\section{Implication dans le forum des partenaires}

Dans le cadre de ce projet, nous avons également participé au forum des partenaires organisé par l'ENSEIRB-MATMECA. Ce forum a été l'occasion de présenter notre travail à un public plus large, composé d'autres étudiants, de professeurs et de professionnels. Nous avons pu échanger avec les participants sur les différentes problématiques liées à la mobilité étudiante et recueillir des retours constructifs sur notre projet.

Lors de notre présentation, nous avons pu effectuer une démonstration de notre application et expliquer les différentes fonctionnalités que nous avons développées. Nous avons reçu des questions pertinentes qui nous ont permis de réfléchir à de nouvelles améliorations pour notre projet, notamment :

\begin{itemize}
\item \textbf{Étendre les fonctionnalités} : ajouter des options supplémentaires pour rendre l'application plus complète, comme la possibilité de suivre l'avancement de la mobilité ou de partager des expériences avec d'autres étudiants.

\item \textbf{Ajouter une fonctionnalité de compétition} : permettre aux étudiants/utilisateurs de se mesurer les uns aux autres en fonction de leurs expériences de mobilité, avec des classements hebdomadaires.
\item \textbf{Ajouter une fonctionnalité de gamification} : rendre l'application plus ludique en intégrant des éléments de jeu, tels que des badges ou des récompenses pour les utilisateurs qui atteignent certains objectifs.
\item \textbf{Élargir notre base d'utilisateurs} : ouvrir notre application à d'autres profils, comme les enseignants ou les entreprises, pour permettre à tout le monde de devenir plus responsable vis-à-vis de ses mobilités.
\item \textbf{Commercialiser notre application} : envisager une commercialisation pour rendre l'outil accessible à un plus grand nombre d'utilisateurs.
\end{itemize}

En effet, le retour le plus récurrent était que notre application est très intéressante et possède un potentiel de développement important. De nombreuses entreprises pourraient notamment être intéressées par une telle solution pour aider leurs employés à gérer leurs mobilités professionnelles.
Pour cela, nous avons également reçu des conseils sur la manière de structurer notre projet pour une éventuelle commercialisation, notamment en termes de marketing et de stratégie de développement. Nous avons aussi discuté de la possibilité de collaborer avec InPulse ou la French Tech afin de lancer la commercialisation de notre application.